\section{Method}

\begin{figure}[t]
\centering
\resizebox{\linewidth}{!}{%
\begin{tikzpicture}[
  node distance=1.45cm and 1.9cm,
  block/.style={draw, rounded corners, align=center, minimum width=3.0cm, minimum height=0.95cm, fill=gray!8},
  knowledge/.style={draw, rounded corners, align=center, minimum width=3.0cm, minimum height=0.9cm, fill=blue!7},
  model/.style={draw, rounded corners, align=center, minimum width=3.0cm, minimum height=0.95cm, fill=green!8},
  output/.style={draw, rounded corners, align=center, minimum width=3.0cm, minimum height=0.95cm, fill=orange!10},
  arrow/.style={-{Latex[length=2mm]}, thick}
]
\node[block] (q) {Medical question\\$q$};
\node[block, right=of q] (first) {First-stage retrieval\\BM25 + dense + RRF};
\node[block, right=of first] (pool) {Top-$M$ candidate\\evidence passages};
\node[knowledge, below=of first] (kg) {Medical constraints\\MeSH hierarchy\\entities\\PrimeKG auxiliary relations};
\node[knowledge, below=of pool] (hg) {Local hypergraph\\question, passage, document,\\entity, concept nodes};
\node[model, right=of pool] (ltr) {Query-group\\LambdaMART reranker};
\node[output, right=of ltr] (out) {Ranked evidence\\top-$k$ passages\\optional paths};
\draw[arrow] (q) -- (first);
\draw[arrow] (first) -- (pool);
\draw[arrow] (pool) -- (ltr);
\draw[arrow] (ltr) -- (out);
\draw[arrow] (kg) -- (hg);
\draw[arrow] (pool) -- (hg);
\draw[arrow] (q) |- (hg);
\draw[arrow] (hg) -- node[below, align=center] {diffusion, centrality,\\hierarchy features} (ltr);
\end{tikzpicture}}
\caption{Overview of KCH-MedRank. The method first builds a high-recall candidate pool, constructs a local knowledge-constrained hypergraph around each question and candidate set, and uses hypergraph and biomedical semantic features in query-group learning-to-rank.}
\label{fig:method-overview}
\end{figure}


\subsection{First-Stage Retrieval}

The first stage retrieves top-$M$ evidence candidates using BM25, dense biomedical retrieval, and hybrid Reciprocal Rank Fusion. The enhanced candidate-generation path adds fielded BM25 with MeSH expansion and MedCPT dual-encoder dense retrieval. Fusion weights are selected on the validation split, where the selected w122 configuration maximizes Recall@100 rather than test performance.

\subsection{Knowledge-Constrained Local Hypergraph}

For each question and top-$M$ candidates, KCH-MedRank builds a local hypergraph with question, passage, document, biomedical entity, MeSH concept, and optional relation nodes. Hyperedges connect query concepts, passage concepts, document-MeSH labels, shared concepts, hierarchy relations, and auxiliary PrimeKG-style relation signals.

The local design avoids constructing a global graph over millions of biomedical passages. It also keeps the method interpretable: feature inspection and case studies can report which concept overlaps, hierarchy links, shared candidate clusters, or hypergraph centrality signals influenced reranking.

\subsection{Learning-to-Rank}

The main supervised reranker is a query-group learning-to-rank model using LightGBM/LambdaMART-style gradient boosted decision trees \citep{ke2017lightgbm}. Features include initial rank scores, BM25 and dense scores, biomedical semantic score, MeSH hierarchy-aware features, entity overlap, local hypergraph diffusion and degree centrality, and auxiliary relation features. Hypergraph diffusion is used as an interpretable feature, not as the sole hand-tuned scoring rule.

\begin{table}[t]
\centering
\small
\caption{Top KCH-MedRank feature importances from the enhanced BioASQ experiment. Interaction features are marked with $\dagger$.}
\label{tab:feature-importance-summary}
\begin{tabular}{lr}
\toprule
Feature & Importance \\
\midrule
Biomedical semantic score & 455 \\
Hybrid score & 404 \\
Dense score & 364 \\
BM25 rank score & 351 \\
Biomedical semantic rank score & 317 \\
Dense rank score & 220 \\
Passage entity count & 219 \\
BM25 score & 196 \\
Hypergraph$\times$inv-rank $\dagger$ & 181 \\
Local num nodes & 157 \\
Hypergraph degree centrality & 145 \\
Shared MeSH term cluster size & 143 \\
Shared MeSH parent cluster size & 131 \\
Mesh cluster$\times$inv-semantic $\dagger$ & 126 \\
Passage MeSH specificity & 118 \\
\bottomrule
\end{tabular}
\end{table}


Feature importance in Table~\ref{tab:feature-importance-summary} shows that semantic and first-stage retrieval signals remain central, while MeSH cluster, passage specificity, and hypergraph centrality features contribute interpretable biomedical and structural evidence. This supports a cautious claim: KCH-MedRank improves reranking by combining retrieval strength, biomedical semantics, and local knowledge-constrained structure, not by relying on one isolated knowledge source.
